\documentclass[UTF8,11pt,a4paper]{ctexart}
\usepackage[margin=24mm]{geometry}
\usepackage{amsmath,amssymb,amsthm,mathtools,bm}
\usepackage{tikz-cd,enumitem,longtable,booktabs}
\usepackage{xurl}
\usepackage[unicode,hidelinks]{hyperref}
\allowdisplaybreaks
\setlength{\emergencystretch}{3em}
\setlength{\parindent}{0pt}
\setlength{\parskip}{0.45em}
\newtheorem*{theorem}{Theorem}
\newtheorem*{lemma}{Lemma}
\newtheorem*{proposition}{Proposition}
\newtheorem*{corollary}{Corollary}
\theoremstyle{definition}\newtheorem*{definition}{Definition}
\theoremstyle{remark}\newtheorem*{remark}{Remark}
\DeclareMathOperator{\Hom}{Hom}
\DeclareMathOperator{\End}{End}
\DeclareMathOperator{\Tr}{Tr}
\DeclareMathOperator{\rank}{rank}
\DeclareMathOperator{\Ind}{Ind}
\DeclareMathOperator{\Res}{Res}
\DeclareMathOperator{\im}{Im}
\DeclareMathOperator{\id}{id}
\newcommand{\R}{\mathbb R}
\newcommand{\C}{\mathbb C}
\newcommand{\Q}{\mathbb Q}
\newcommand{\Z}{\mathbb Z}
\newcommand{\N}{\mathbb N}
\newcommand{\E}{\mathbb E}
\newcommand{\Prob}{\mathbb P}
\newcommand{\eps}{\varepsilon}
\begin{document}
\section*{091｜Brunn--Minkowski体积下界与积分不等式的张量化}
\textbf{作者／集体来源：}Terence Tao

\textbf{作品：}The Brunn--Minkowski inequality for nilpotent groups，2011-09-16

\textbf{原文入口：}\url{https://terrytao.wordpress.com/2011/09/16/the-brunn-minkowski-inequality-for-nilpotent-groups/}

\textbf{原文位置与计题范围：}欧氏部分：式(1)、Lemma 1、Theorem 2及其后完整论证；仅计欧氏Brunn--Minkowski问题一道。

\textbf{获取日期：}2026-09-28。\quad\textbf{复用依据：}作者About页面明确允许带来源复制合理篇幅（例如单篇文章）；本题取其中欧氏部分。

\textbf{整理归属：}下列题面与证明取自所列人类来源，按注明范围转录、摘编；LaTeX环境、标题、换行及中文说明由助手整理。编辑调整在题内说明。

\section*{作者命题与人类证明}

\paragraph{Brunn--Minkowski inequality (equation (1) in the source).}
One of the fundamental inequalities in convex geometry is the Brunn--Minkowski inequality, which asserts that if $A,B$ are two non-empty bounded open subsets of $\mathbf R^d$, then
\begin{equation}
\mu(A+B)^{1/d}\geq\mu(A)^{1/d}+\mu(B)^{1/d},\tag{1}
\end{equation}
where
\[
A+B:=\{a+b:a\in A,b\in B\}
\]
is the sumset of $A$ and $B$, and $\mu$ denotes Lebesgue measure. The estimate is sharp, as can be seen by considering the case when $A,B$ are convex bodies that are dilates of each other, thus $A=\lambda B:=\{\lambda b:b\in B\}$ for some $\lambda>0$, since in this case one has $\mu(A)=\lambda^d\mu(B)$, $A+B=(\lambda+1)B$, and $\mu(A+B)=(\lambda+1)^d\mu(B)$.

\paragraph{Lemma 1 (One-dimensional Brunn--Minkowski).}
If $A,B,C\subset\mathbf R$ are non-empty measurable sets with $A+B\subset C\subset\mathbf R$, then
\[\mu(C)\geq\mu(A)+\mu(B).\]
\begin{proof}
By inner regularity we may assume that $A,B$ are compact. The claim then follows since $C$ contains the sets $\sup(A)+B$ and $A+\inf(B)$, which meet only at a single point $\sup(A)+\inf(B)$.
\end{proof}
For the higher dimensional case, the inequality can be established from the Pr\'ekopa--Leindler inequality:

\paragraph{Theorem 2 (Pr\'ekopa--Leindler inequality in $\mathbf R^d$).}
Let $0<\theta<1$, and let $f,g,h:\mathbf R^d\to\mathbf R$ be non-negative measurable functions obeying the inequality
\begin{equation}
h(x+y)\geq f(x)^{1-\theta}g(y)^\theta\tag{2}
\end{equation}
for all $x,y\in\mathbf R^d$. Then we have
\begin{equation}
\int_{\mathbf R^d}h\geq
\frac{1}{(1-\theta)^{d(1-\theta)}\theta^{d\theta}}
\left(\int_{\mathbf R^d}f\right)^{1-\theta}
\left(\int_{\mathbf R^d}g\right)^\theta.\tag{3}
\end{equation}
This inequality is usually stated using $h((1-\theta)x+\theta y)$ instead of $h(x+y)$ in order to eliminate the ungainly factor $\frac{1}{(1-\theta)^{d(1-\theta)}\theta^{d\theta}}$. However, we formulate the inequality in this fashion in order to avoid any reference to the dilation maps $x\mapsto\lambda x$; the reason for this will become clearer later.

The Pr\'ekopa--Leindler inequality quickly implies the Brunn--Minkowski inequality. Indeed, if we apply it to the indicator functions $f:=1_A,g:=1_B,h:=1_{A+B}$ (which certainly obey (2)), then (3) gives
\[
\mu(A+B)^{1/d}\geq
\frac{1}{(1-\theta)^{1-\theta}\theta^\theta}
\mu(A)^{\frac{1-\theta}{d}}\mu(B)^{\frac{\theta}{d}}
\]
for any $0<\theta<1$. We can now optimise in $\theta$; the optimal value turns out to be
\[
\theta:=\frac{\mu(B)^{1/d}}{\mu(A)^{1/d}+\mu(B)^{1/d}}
\]
which yields (1).

To prove the Pr\'ekopa--Leindler inequality, we first observe that the inequality tensorises in the sense that if it is true in dimensions $d_1$ and $d_2$, then it is automatically true in dimension $d_1+d_2$. Indeed, if $f,g,h:\mathbf R^{d_1}\times\mathbf R^{d_2}\to\mathbf R^+$ are measurable functions obeying (2) in dimension $d_1+d_2$, then for any $x_1,y_1\in\mathbf R^{d_1}$, the functions $f(x_1,\cdot),g(y_1,\cdot),h(x_1+y_1,\cdot):\mathbf R^{d_2}\to\mathbf R^+$ obey (2) in dimension $d_2$. Applying the Pr\'ekopa--Leindler inequality in dimension $d_2$, we conclude that
\[
H(x_1+y_1)\geq
\frac{1}{(1-\theta)^{d_2(1-\theta)}\theta^{d_2\theta}}
F(x_1)^{1-\theta}G(y_1)^\theta
\]
for all $x_1,y_1\in\mathbf R^{d_1}$, where $F(x_1):=\int_{\mathbf R^{d_2}}f(x_1,x_2)\,dx_2$ and similarly for $G,H$. But then if we apply the Pr\'ekopa--Leindler inequality again, this time in dimension $d_1$ and to the functions $F$, $G$, and $(1-\theta)^{d_2(1-\theta)}\theta^{d_2\theta}H$, and then use the Fubini--Tonelli theorem, we obtain (3).

From tensorisation, we see that to prove the Pr\'ekopa--Leindler inequality it suffices to do so in the one-dimensional case. We can derive this from Lemma 1 by reversing the ``Pr\'ekopa--Leindler implies Brunn--Minkowski'' argument given earlier, as follows. We can normalise $f,g$ to have sup norm $1$. If (2) holds (in one dimension), then the super-level sets $\{f>\lambda\},\{g>\lambda\},\{h>\lambda\}$ are related by the set-theoretic inclusion
\[
\{h>\lambda\}\supset\{f>\lambda\}+\{g>\lambda\}
\]
and thus by Lemma 1
\[
\mu(\{h>\lambda\})\geq\mu(\{f>\lambda\})+\mu(\{g>\lambda\})
\]
whenever $\lambda\leq1$. On the other hand, from the Fubini--Tonelli theorem one has the distributional identity
\[
\int_{\mathbf R}h=\int_0^\infty\mu(\{h>\lambda\})\,d\lambda
\]
(and similarly for $f,g$, but with $\lambda$ restricted to $(0,1)$), and thus
\[
\int_{\mathbf R}h\geq\int_{\mathbf R}f+\int_{\mathbf R}g.
\]
The claim then follows from the weighted arithmetic mean-geometric mean inequality $(1-\theta)x+\theta y\geq x^{1-\theta}y^\theta$.

\section*{整理说明与可修改标注（助手）}
\textbf{标准先修：}Lebesgue测度的内正则性、Fubini--Tonelli、加权算术—几何平均不等式。

\textbf{本题仍需完成的工作：}将集合加法改写为超水平集包含，再通过积分不等式的张量化从一维推进到任意维，最后优化参数恢复精确体积幂次。

\textbf{取材说明：}保留原文论证顺序及式(1)--(3)。省略与本题无关的等周应用、后半篇幂零Lie群推广和评论。原文幂零推广留有练习，本题不计该推广。

\textbf{$c$：}欧氏集合和的体积估计

\textbf{$a$：}示性函数；超水平集；层蛋糕积分；Fubini--Tonelli；积分不等式张量化；参数优化。

\texttt{cross\_theory\_status = yes}

\textbf{具体联系及位置：}Theorem 2后两段把集合加法转为函数乘法下界，通过切片积分把不同维数接合；其前的参数优化段返回集合体积。

这些标注是非穷尽初稿；未列出不表示未使用。
\end{document}
